\documentclass[12pt,a4paper,two side]{article}
\usepackage[left=2.5cm,right=2.5cm,top=3cm,bottom=3cm]{geometry}
\usepackage[utf8]{vietnam}
\usepackage{amsmath,amsfonts,amssymb,amsthm}
\usepackage{pgf,tikz}
\usepackage{mathrsfs}
\usetikzlibrary{arrows}
\usepackage{enumerate}
\usepackage{tcolorbox}
\usepackage{multicol}
\usepackage{graphicx}
\usepackage{wrapfig}
\usepackage{float}
\usepackage{fancyhdr}
\pagestyle{fancy}
\fancyhf{}
\rhead{\thepage}
\lhead{\textbf{\textbf{Câu lạc bộ Toán A1}} - \textbf{Hotline:}  \ \textsc{035 290 3286\ -\ 034 761 1986}}
\fancypagestyle{footer}{\fancyhf{}\renewcommand{\headrulewidth}{0pt}\renewcommand{\footrulewidth}{0.5pt}\fancyfoot[L]{\textbf{Hotline:} \ \textsc{035 290 3286\ -\ 034 761 1986}}\fancyfoot[R]{\thepage}}

\usepackage{background}
\backgroundsetup{
	scale=1,
	color=black,
	opacity=0.08,
	angle=0,
	contents={%
		\includegraphics[scale=0.7]{ToKien.jpg}
	}%
}

\setlength{\parindent}{0pt}
\theoremstyle{definition}
\newtheorem{prob}{Bài}
\newtheorem{lemma}{Bổ đề}
\begin{document}
	\large
	\thispagestyle{footer}
	\begin{minipage}{0.3\textwidth}
		\begin{center}
			\includegraphics[scale=0.12]{ToKien.jpg}
		\end{center}
	\end{minipage}
	\begin{minipage}{0.7\textwidth}
		\begin{tcolorbox}[colback=red!5,colframe=red!75!black]
			\begin{center}
				{\large \textbf{ĐÁP ÁN BÀI TẬP VỀ NHÀ BUỔI 21}} 
			\end{center}
			\tcblower
			{\large \textbf{Trợ giảng:} Ngô Văn Đăng}
			\flushleft{\large \textbf{Lớp:} 9H4  \hfill \textbf{\textit{Ngày 10/11/2025}}}
		\end{tcolorbox}
	\end{minipage}
	\vspace{0.5cm}
\begin{prob}
Cho tứ giác $ABCD$ nội tiếp đường tròn $(O)$. Gọi $E$ là giao điểm của $AC$ và $BD$. $F$ là giao điểm thứ hai của hai đường tròn $(EAD)$ và $(EBC)$.
\begin{enumerate}[a)]
	\item Chứng minh tứ giác $ABOF$ nội tiếp.
	\item Chứng minh $\widehat{OEF} = 90^{\circ}$.
\end{enumerate}
\end{prob}
\textit{\textbf{Lời giải: }}
\begin{figure}[H]
	\centering
	\includegraphics[scale=0.7]{bài1}
\end{figure}
a) Ta có: $\widehat{ADE}=\widehat{AFE}$ và $\widehat{EFB}=\widehat{ACB}$\\
Nên $\widehat{AFE}=\widehat{AFB}=\widehat{AFE}+\widehat{EFB}=2\widehat{ADB}=\widehat{AOB}$\\
Suy ra $ABOF$ là tứ giác nội tiếp.\\
b) Ta có: $\widehat{OFB}=\widehat{OAB}=90^\circ -\dfrac{1}{2}\widehat{AOB}=90^\circ - \widehat{ACB}$\\
và $\widehat{EFB}=\widehat{ECB}$\\
Suy ra $\widehat{OFE}=\widehat{OFB}+\widehat{BFE}=90^\circ$\\

\begin{prob}
Cho tam giác nhọn $ABC$ nội tiếp $(O)$, với $AB < AC$. Gọi $AD$ là phân giác góc $A$, cắt $(O)$ tại $K$. Đường thẳng $AC$ cắt đường tròn $(ABD)$ tại $L$ (khác $A$).
\begin{enumerate}[a)]
	\item Chứng minh rằng $BL \parallel CK$.
	\item Qua $L$ kẻ đường thẳng vuông góc với $BC$, cắt $(ABD)$ tại $P$. Chứng minh rằng $B, P, O$ thẳng hàng.
	\item Đường thẳng $AP$ cắt $(O)$ tại $Q$. Chứng minh rằng $CQ \perp BK$.
\end{enumerate}
\end{prob}
\textit{\textbf{Lời giải: }}
\begin{figure}[H]
	\centering
	\includegraphics[scale=0.7]{bài2}
\end{figure}
a) Ta có: $\widehat{ALB}=\widehat{ADB}=\dfrac{1}{2}(sdAB+sdKC)=\dfrac{1}{2}(sdAB+sdBK)=\widehat{ACK}$\\
Suy ra $BL \parallel CK$.\\
b) Ta có: $\widehat{PBC}=90^\circ-\widehat{BPL}=90^\circ-\widehat{BAC}=\widehat{OBC}$\\
Suy ra $O,P,B$ thẳng hàng.\\
c) Gọi $CQ$ cắt $BK$ tại $T$.\\
Ta có: $\widehat{TQB}+\widehat{TBQ}=\widehat{BAC}+\widehat{DAP}=\dfrac{1}{2}\widehat{BOC}+\widehat{DBO}=90^\circ$\\

\begin{prob}
Cho tam giác $ABC$ nội tiếp đường tròn $(O)$. Gọi $I$ là tâm nội tiếp. Các tia phân giác $AI, BI, CI$ cắt $(O)$ lần lượt tại $M, N, P$ (khác các đỉnh tương ứng).
\begin{enumerate}[a)]
	\item Chứng minh $MB = MC = MI$.
	\item Chứng minh $I$ là trực tâm tam giác $MNP$.
	\item Các đường thẳng $MN, MP$ lần lượt cắt $AC, AB$ tại $E, D$. Chứng minh $E, I, D$ thẳng hàng.
\end{enumerate}
\end{prob}
\textit{\textbf{Lời giải: }}
\begin{figure}[H]
	\centering
	\includegraphics[scale=0.6]{bài3}
\end{figure}
a) Ta có: $\widehat{MIB}=\dfrac{1}{2}(sdAN+sdBM)=\dfrac{1}{2}(sdNC+sdMC)=\widehat{MBI}$\\
Suy ra $\triangle MIB$ cân tại $M$ nên $MB=MI$ \\
Mà có $MB=MC$ nên có $MI=MB=MC$\\
b)
Tương tự ta có $PI=PB$ nên $PI=PB=PA$\\
Tương tự cho điểm $N$ ta có: $NA=NC=NI$\\
Suy ra $NP$ là trung trực của $AI$\\
Tương tự suy ra $I$ là trực tâm tam giác $MNP$\\
c) Theo kết quả quen thuộc ta có $PD.PM=PB^2$ nên $PM.PD=PI^2$\\
Suy ra $\widehat{PID}=\widehat{PMA}=\widehat{PCB}$\\
Suy ra $ID \parallel BC$\\
Tương tự ta cũng chứng minh được $IE \parallel BC$\\
Suy ra $I,D,E$ thẳng hàng.\\



\begin{prob}
Cho tam giác $ABC$ nội tiếp đường tròn $(O)$. Điểm $E$ thuộc cung nhỏ $AC$. Đường thẳng $AE$ cắt các tiếp tuyến tại $B$ và $C$ của $(O)$ lần lượt tại $P, Q$. Đường thẳng $BQ$ cắt $CP$ tại $F$.
\begin{enumerate}[a)]
	\item Chứng minh $AC \parallel BP$.
	\item Chứng minh $\triangle BCP \sim \triangle CQB$.
	\item Chứng minh tứ giác $BPEF$ nội tiếp.
	\item Chứng minh $EF$ đi qua trung điểm của $BC$.
\end{enumerate}
\end{prob}
\textit{\textbf{Lời giải: }}
\begin{figure}[H]
	\centering
	\includegraphics[scale=0.6]{bài4}
\end{figure}
a) Ta có: $\widehat{PBA}=\widehat{BAC}=60^\circ$ nên $AC \parallel BP$\\
b) Tương tự cũng có $AB \parallel QC$\\
Suy ra $\widehat{APB}=\widehat{CAQ}$ và $\widehat{PAB}=\widehat{AQC}$\\
Suy ra $\triangle PAB \backsim \triangle AQC (g.g) \Rightarrow \dfrac{AB}{QC}=\dfrac{BP}{AC} \Rightarrow \dfrac{BC}{QC}=\dfrac{BP}{BC}$ và $ \widehat{PBC}=\widehat{BCQ}=120\circ$\\
Suy ra $\triangle BCP \sim \triangle CQB(c.g.c)$
c) Ta có: $\widehat{BFP}=\widehat{FBC}+\widehat{FCB}=\widehat{FBC}+\widehat{CQB}=60^\circ$\\
Và $\widehat{BFP}=\widehat{ACB}=60^\circ$\\
Suy ra $PEFB$ là tứ giác nội tiếp.\\
d) Gọi $EF$ cắt $BC$ tại $M$\\
Ta có: $\widehat{FBM}=\widehat{BPC}=\widehat{BEM}$\\
Suy ra $\triangle BMF \backsim \triangle EMB (g.g)$\\
Suy ra $ MB^2=MF.ME$\\
Tương tự ta cũng chứng minh được $MC^2=ME.MF$\\
Suy ra $MC=MB$ (đpcm)\\


\begin{prob}
Cho hình thang $ABCD$ với $AD, BC$ là hai cạnh đáy; $BC > AD$, $BC = BD = 1$, $AB = AC, CD < 1$, $\widehat{BAC} + \widehat{BDC} = 180^{\circ}$. Gọi $E$ là điểm đối xứng với $D$ qua $BC$.
\begin{enumerate}[a)]
	\item Chứng minh bốn điểm $A, C, E, B$ cùng nằm trên một đường tròn và $\widehat{BEC} = 2\widehat{AEC}$.
	\item Đường thẳng $AB$ cắt đường thẳng $CD$ tại $K$, đường thẳng $BC$ cắt đường thẳng $AE$ tại $F$. Chứng minh $FA = FD$ và đường thẳng $FD$ tiếp xúc với đường tròn $(ADK)$.
	\item Tính độ dài cạnh $CD$.
\end{enumerate}
\end{prob}
\textit{\textbf{Lời giải: }}
\begin{figure}[H]
	\centering
	\includegraphics[scale=0.5]{bài5}
\end{figure}
a) Ta có: $\widehat{BAC}+\widehat{BEC}= \widehat{BAC}+\widehat{BDC}=180^{\circ}$\\
Suy ra $A,E,C,B$ cùng nằm trên một đường tròn.\\
Ta có: $\widehat{BEC}=\widehat{AEC}+\widehat{AEB}=\widehat{ACB}+\widehat{ABC}=2\widehat{ABC}=2\widehat{AEC}$\\
b) Ta có: $\widehat{FAD}=\widehat{AFB}=\widehat{CFE}=\widehat{CFD}=\widehat{FDA}$\\
Suy ra tam giác $FAD$ cân tại $F$ nên $FA=FD$\\
Ta có: $\widehat{BKC}=180^\circ -\widehat{KBC}-\widehat{KCB}=180^\circ -\widehat{ABC}-\widehat{ECB}$\\
$=\dfrac{1}{2}(sdAB+sdCE)=\widehat{AFB}=\widehat{FAD}$\\
Suy ra $FA$ là tiếp tuyến của $(ADK)$\\
Tương tự cũng suy ra được $FD$ là tiếp tuyến của $(ADK)$\\
c)\textit{(Lời giải của bạn Lê Vũ Khánh Linh $9H4$ $ 25-26$ )} \\
Áp dụng định lí $Ptolemy$ cho tứ giác nội tiếp $ACEB$\\
Ta có: $AE.BC=AC.BE+AB.CE$\\
Suy ra $ 2AF=AC(1+CE)$\\
Suy ra $\dfrac{2}{1+CE}=\dfrac{AC}{AF}=\dfrac{BE}{BF}=\dfrac{BC}{BF}=\dfrac{BF+FC}{BF}=1+\dfrac{FC}{BF}$\\
Mà ta có $EF$ là phân giác $\widehat{BEC}$ nên $\dfrac{CF}{BF}=\dfrac{CE}{BE}=CE$\\
Suy ra $\dfrac{2}{1+CE}=1+CE$\\
Suy ra $ (CE+1)^2=2 \Rightarrow CE+1= \sqrt{2}$\\
Suy ra $CE=\sqrt{2}-1=CD$\\

\end{document}